% 出席確認クイズ 第07回　ケーススタディ：企業のAI活用
%   attendance/attendance07.html から生成（解説は本ガイド用に加筆）。

\quizq{1}{AI導入の効果を測る際に、まず必要なことはどれでしょうか?}%
  {効果測定は行わない}%
  {導入費用が高いほど効果が高いとみなす}%
  {\correct{導入前に目的とKPI(評価指標)を決めておく}}%
  {導入後になんとなく満足度を聞く}%
  {正解は (c)。効果を測るには、導入前に目的と評価指標（KPI）を定め、導入前後で比較できるようにしておく必要があります。(d) のように後からなんとなく満足度を聞いても比較できません。(b) 費用と効果は別のものです。}

\quizq{2}{「PoC(概念実証)」の説明として適切なものはどれでしょうか?}%
  {契約書の一種}%
  {本番システムの名称}%
  {\correct{本格導入の前に小規模で実現可能性や効果を検証すること}}%
  {AIの学習アルゴリズムの名前}%
  {正解は (c)。PoC（Proof of Concept、概念実証）は、本格導入の前に小規模で技術的な実現可能性や効果を確かめる取り組みです。PoC が次に進まず繰り返される「PoC止まり」が多くの企業で課題になっており、PoC の段階で本番展開の条件を決めておくことが重要です。}

\quizq{3}{AI導入が失敗しやすい典型的な要因として適切なものはどれでしょうか?}%
  {データの整備をしすぎる}%
  {現場を巻き込みすぎる}%
  {\correct{目的が曖昧なまま技術先行で導入する}}%
  {目的が明確すぎる}%
  {正解は (c)。「AIを使うこと」自体が目的化し、解決したい課題が曖昧なまま技術先行で導入するのが典型的な失敗要因です。(a)(b) は「やりすぎ」を装った選択肢ですが、データの整備と現場の巻き込みはむしろ成功要因です。}

\quizq{4}{顧客対応でのAI活用例として一般的なものはどれでしょうか?}%
  {商品の物理的な配送}%
  {電力の供給}%
  {店舗の清掃}%
  {\correct{チャットボットによる問い合わせ対応}}%
  {正解は (d)。問い合わせ対応のチャットボットは、顧客対応における AI 活用の代表的な例です。(a)(b)(c) は物理的な作業で、生成AIが直接担う領域ではありません。チャットボットでも、回答の正確さと人への引き継ぎの設計が成否を分けます。}

\quizq{5}{他社の事例を読み解く際の姿勢として適切なものはどれでしょうか?}%
  {有名企業の事例だけを見ればよい}%
  {成功事例をそのまま真似すれば必ず成功する}%
  {失敗事例には見る価値がない}%
  {\correct{業種・規模・データの条件が自社と異なる点を踏まえて比較する}}%
  {正解は (d)。事例は業種・規模・データの条件などの前提が異なるため、自社との違いを踏まえて何が応用できるかを考えることが大切です。(b) 成功事例の模倣が必ず成功するわけではなく、(c) 失敗事例からはつまずきの要因を学べます。(a) 規模の近い企業の事例の方が参考になることも多いです。}
